\documentclass[11pt]{article}

% --- Encoding & fonts ---
\usepackage[utf8]{inputenc}
\usepackage[T1]{fontenc}
\usepackage{lmodern}

% --- Layout ---
\usepackage[margin=1in]{geometry}
\usepackage{microtype}
\usepackage{setspace}
\onehalfspacing
\usepackage{booktabs}
\usepackage{enumitem}

% --- Citations & links (natbib + bibtex; compiles on vanilla Overleaf with pdfLaTeX) ---
\usepackage[dvipsnames]{xcolor}
\usepackage[round]{natbib}
\bibliographystyle{plainnat}
\usepackage[colorlinks=true,linkcolor=black,citecolor=NavyBlue,urlcolor=NavyBlue]{hyperref}

\title{\textbf{Tokenizer-Free Embeddings for Low-Resource Question Answering:\\[2pt]
Byte-Level versus Subword Multilingual Retrieval}}
\author{Aarush Vinod\\ \small Sorva Health \\ \small \texttt{avinod@sorvahealth.com}}
\date{\today}

\begin{document}
\maketitle

% =====================================================================
\begin{abstract}
\noindent
Answering questions in low-resource languages increasingly relies on retrieval-augmented
generation (RAG), where a retriever finds the passage containing the answer. Yet the multilingual
embeddings these retrievers use rely on English-centric subword tokenizers. These tokenizers segment
low-resource languages into far more tokens than they do for high-resource languages. We ask whether
tokenizer-free, byte-level embeddings yield better retrieval. To test this, we distill SONAR
embeddings into byte-level (ByT5) and subword (mT5) students across 3 model sizes, with fixed training
objectives and parameters. We evaluate both embeddings on passage retrieval (Belebele) and deep QA
tasks (MIRACL, AfriQA, IndicQA), for six low-resource languages (Telugu, Tamil, Marathi, Amharic,
Kinyarwanda, Hausa). Preliminary results show that byte embeddings perform better on passage
retrieval, although the advantage diminishes as parameter count increases. For deep retrieval, byte
embeddings perform significantly better, with the smallest byte-level model surpassing the largest
subword model. This project aims to show that byte-level embeddings yield fairer multilingual
embeddings well-suited for low-resource languages.
\end{abstract}

% =====================================================================
\section{Introduction}
\label{sec:intro}

Large language models are increasingly the interface through which people seek information, and for
the many languages with little training data this is made workable by retrieval-augmented generation
(RAG): rather than answer from parametric memory alone, the system first \emph{retrieves} a passage
likely to contain the answer and then conditions its response on that passage \citep{lewis2020rag}.
The quality of the final answer is therefore bounded by the quality of the retriever. For
low-resource languages, however, the multilingual sentence embeddings that power retrieval are built
on subword tokenizers trained on English-dominated corpora, which fragment low-resource and
non-Latin-script languages into far more tokens than high-resource ones \citep{petrov2023, ahia2023}.
This ``tokenization tax'' degrades representation quality for exactly the languages that are least
served today and that stand to benefit most from better retrieval.

This proposal asks whether removing the tokenizer helps. We compare \emph{byte-level} (ByT5) and
\emph{subword} (mT5) multilingual sentence embeddings, distilled from a shared frozen SONAR teacher
under a fixed recipe so that only the input representation varies, and evaluate them as retrievers
for six low-resource languages spanning four scripts and five language families. Concretely, we
investigate four research questions:

\begin{description}[leftmargin=2.2em,style=nextline,itemsep=2pt]
  \item[\textbf{RQ1.}] Do byte-level embeddings retrieve better than subword embeddings for
  low-resource languages, on both shallow passage retrieval and deep open-domain (QA) retrieval?
  \item[\textbf{RQ2.}] How does any advantage depend on model size --- in particular, can the
  smallest byte-level model match or surpass the largest subword model?
  \item[\textbf{RQ3.}] Is a byte advantage attributable to the \emph{representation} rather than to
  raw parameter count, when transformer capacity is held equal across the two tokenizers?
  \item[\textbf{RQ4.}] Can an intrinsic, downstream-free measure of tokenization (such as
  cross-lingual \emph{parity}) characterize the inequity that motivates the approach, and how do byte
  and subword tokenizers compare on it?
\end{description}

\noindent
Preliminary results suggest the answer to RQ1 is yes, and most strongly for deep retrieval, where the
smallest byte model already surpasses the largest subword model. Section~\ref{sec:related} situates
the work in prior literature on the tokenization tax and tokenizer-free modeling, and
Section~\ref{sec:proposed} sketches the planned study.

% =====================================================================
\section{Related Work}
\label{sec:related}

Modern NLP serves the world's languages unequally, and a growing body of work locates one root of
that inequality in the tokenizer. \citet{joshi2020} document a steep resource hierarchy in which most
languages, including the majority spoken across the Global South, receive little data, tooling, or
evaluation. Subword tokenizers, trained on corpora dominated by high-resource languages, inherit and
amplify this imbalance. \citet{petrov2023} show that identical content translated across languages is
segmented into very different numbers of tokens, with disparities reaching $15\times$ and non-Latin
scripts among the worst affected. \citet{ahia2023} quantify the effect on parallel FLORES-200 text,
where mid-resourced languages with their own scripts such as Telugu and Georgian require up to
$5\times$ more tokens than English. Because commercial models bill per token, this tokenization tax is
also economic: speakers of those languages pay several times more for equivalent content.

The penalty is structural rather than incidental. \citet{somide2026} finds that every African
language studied carries a per-language premium over English, and that even the best available subword
tokenizer narrows but never closes the gap. \citet{lundin2025tokentax} report that higher subword
fertility, the number of tokens per word, is associated with lower downstream accuracy on
African-language QA, though the relationship is correlational and partly confounded with
pretraining-data volume. Both are recent preprints, and their exact multipliers should be read with
that caveat, and with the tokenizer version on which they were measured, since newer vocabularies have
reduced the worst cases.

The natural response is to remove the tokenizer. \citet{clark2022canine} argue that fixed subword
vocabularies are not equally suited to all languages and operate directly on characters, discarding
the vocabulary at inference; CANINE outperforms a larger subword mBERT by $2.8$ F1 on multilingual
TyDi QA while using $28\%$ fewer parameters. ByT5 \citep{xue2022byt5}, the byte-level model we build
on, matches or beats the equally sized subword mT5 \citep{xue2021mt5} at small scale, attributed to
reallocating the parameters that subword models lock in a large vocabulary-embedding matrix into the
transformer itself. Once viewed as a small-model curiosity, byte-level modeling now scales: the Byte
Latent Transformer \citep{pagnoni2024blt} matches a FLOP-controlled Llama~3 up to 8B parameters.
Together these establish both feasibility and a precedent for a smaller tokenizer-free model
outperforming a larger subword one.

The application that most exposes the tokenizer's cost is question answering, which for low-resource
languages increasingly depends on retrieving an answer-bearing passage before generation. Yet existing
evidence for tokenizer-free gains concentrates on classification and extractive QA, with no prior
study isolating a byte advantage for dense retrieval in low-resource languages. This is the gap our
work targets.

Finally, the tokenizer literature offers an intrinsic, downstream-free way to characterize
multilingual fairness. Tokenization parity, the token-length ratio of a language to English on
parallel text \citep{petrov2023, ahia2023}, is computable for any tokenizer, including byte-level ones
where it reduces to average bytes per sentence, making it the one number that compares byte and
subword on equal terms; complements include the Gini coefficient of per-language cost
\citep{lee2026equity} and the information-theoretic R\'enyi efficiency \citep{zouhar2023}. Two cautions
apply. First, byte-level input reduces but does not eliminate cross-lingual disparity, which still
exceeds $4\times$ for some pairs because UTF-8 length varies by script \citep{petrov2023}. Second,
these intrinsic measures do not reliably predict downstream quality, which is task-dependent
\citep{lee2026equity, ali2023tokenizer}. Parity is therefore best reported as a descriptive measure of
multilingual equity, not a forecaster of task performance.

% =====================================================================
\section{Proposed Work}
\label{sec:proposed}

We distill SONAR \citep{duquenne2023sonar} sentence embeddings into byte-level (ByT5) and subword
(mT5) students at three model sizes, holding the training data, supervision, and objective fixed so
that only the input representation differs (addressing RQ1--RQ2). We evaluate retrieval on Belebele
\citep{bandarkar2023belebele} across all languages, and deep open-domain QA retrieval --- MIRACL
\citep{zhang2023miracl}, AfriQA \citep{ogundepo2023afriqa}, and IndicQA \citep{doddapaneni2023indicxtreme}
--- for the six low-resource languages. To separate representation from capacity (RQ3) we add a
parameter-matched ablation in which byte and subword students share an identical transformer and differ
only in the input embedding and tokenizer. To characterize the underlying inequity (RQ4) we report
tokenization parity as an intrinsic, downstream-free fairness measure alongside the retrieval results.
\emph{[Expand with the detailed training protocol, baselines, statistical treatment, and timeline.]}

% =====================================================================
\bibliography{references}

\end{document}
